\section*{Femur case study: robustness on anatomical geometry}
\label{sec:femur-comparison}

\paragraph{Motivation.}
The proximal femur provides a challenging anatomical setting for coupled mechanics--stimulus--density remodeling: the model should simultaneously yield (i) a thin, high-density cortical shell and (ii) a spatially heterogeneous trabecular core under a small set of representative daily loading cases.
In a standard mechanostat formulation a single reference stimulus $\psi_{\mathrm{ref}}$ is used throughout the domain \cite{Weinans1992_BehaviorAdaptiveBoneRemodeling,MartinezReina2016_UniquenessBRM}.
However, experimental and conceptual work suggests that ``set points'' and adaptation thresholds are not universal constants: they may shift with loading history (mechanical usage set-point concept) \cite{Frost1992_SetPointsOsteoporosis}, depend on cellular set-point distributions \cite{Mullender1998_SetPointTrabecularArchitecture}, and can differ between cortical and cancellous compartments \cite{Yang2021_CancellousGreaterThreshold,Schulte2013_LocalMechanicalStimuli}.
This motivates our compartment-aware reference stimulus $\hat{\psi}_{\mathrm{ref}}(\widetilde\rho)$ in \S\ref{subsec:stim-nd}, which blends trabecular and cortical reference values using the same density-dependent transition $w(\widetilde\rho)$ as the stiffness exponent blend.

\paragraph{Ratio sweep definition.}
We parameterize the two reference values by their ratio
\[
r:=\psi_{\mathrm{ref}}^{\mathrm{cort}}/\psi_{\mathrm{ref}}^{\mathrm{trab}},
\]
and, to decouple ``overall'' mechanostat scaling from the cortical--trabecular contrast, we hold the geometric mean $\psi_0:=\sqrt{\psi_{\mathrm{ref}}^{\mathrm{trab}}\psi_{\mathrm{ref}}^{\mathrm{cort}}}$ fixed while sweeping $r$:
\[
\psi_{\mathrm{ref}}^{\mathrm{trab}}(r)=\psi_0/\sqrt{r},\qquad
\psi_{\mathrm{ref}}^{\mathrm{cort}}(r)=\psi_0\sqrt{r}.
\]
All runs use the accelerated \textbf{StiffSetup} femur configuration (\texttt{stiff\_params\_femur.json}) over $T=500$ days, with identical loading cases and all other parameters fixed.

\paragraph{Reference density and mapping.}
As a diagnostic population reference we use the HFValid-derived template femur density (population average), represented as a volumetric CT image and mapped onto the FE mesh using inverse-distance weighting (IDW) interpolation (\texttt{intensity\_mapper.py}).
For a consistent comparison with the bounded density model, the mapped CT field is linearly rescaled to $[\rho_{\min},\rho_{\max}]$ (cf.\ \texttt{analysis/psi\_ref\_ratio\_analysis.py}).
Because the reference is a population template mapped and rescaled to the FE mesh (not subject-specific longitudinal ground truth), the comparisons below are used as diagnostic similarity measures for robustness and relative trends across model variants rather than as clinical validation.

\paragraph{Comparison strategy and metrics.}
We compare the final simulated density fields $\rho_r(\cdot,T)$ to the mapped CT reference $\rho_{\mathrm{CT}}$ in three complementary ways:
\begin{enumerate}
  \item \textbf{Global similarity:} discrete nodal RMSE/MAE over CG1 degrees of freedom,
  \[
  \mathrm{RMSE}(r)=\sqrt{\frac{1}{N}\sum_{j=1}^N\big(\rho_r(x_j,T)-\rho_{\mathrm{CT}}(x_j)\big)^2},\qquad
  \mathrm{MAE}(r)=\frac{1}{N}\sum_{j=1}^N\big|\rho_r(x_j,T)-\rho_{\mathrm{CT}}(x_j)\big|,
  \]
  together with a Pearson correlation as a scale-insensitive summary of spatial pattern similarity.
  \item \textbf{Compartment-restricted similarity:} the same metrics restricted to CT-defined trabecular and cortical masks,
  \[
  \Omega_{\mathrm{trab}}:=\{x:\rho_{\mathrm{CT}}(x)\le \rho_{\mathrm{trab}}^{\max}\},\qquad
  \Omega_{\mathrm{cort}}:=\{x:\rho_{\mathrm{CT}}(x)\ge \rho_{\mathrm{cort}}^{\min}\},
  \]
  using the transition densities from the model (\S\ref{sec:strong-forms-nd}).
  \item \textbf{Feature preservation (contrast):} the cortical--trabecular density contrast
  \(
  C(r):=\bar\rho_{r,\mathrm{cort}}-\bar\rho_{r,\mathrm{trab}},
  \)
  compared to the CT contrast $C_{\mathrm{CT}}$ computed on the same masks.
\end{enumerate}
In addition, we track the \emph{global} density evolution via the domain-mean density $\bar\rho(t)$ reconstructed from the conserved mass integral (logged as \texttt{total\_mass\_g} in \texttt{steps.csv}), which provides a coarse check of transient behavior and settling time.

\paragraph{Hypothesis and test.}
We already observe qualitatively that the special case $\psi_{\mathrm{ref}}^{\mathrm{trab}}=\psi_{\mathrm{ref}}^{\mathrm{cort}}$ (i.e.\ $r=1$) suppresses key spatial features in the simulated density distribution.
We therefore hypothesize that a non-unity ratio $r\neq 1$ is required to recover anatomically plausible cortical--trabecular contrast and to improve similarity to the population-template reference under the same loads.
We test this hypothesis by (i) sweeping $r$ at fixed $\psi_0$, (ii) ranking runs by similarity-to-template metrics and contrast preservation, and (iii) visually confirming differences via distribution plots and (optionally) spatial residuum fields in ParaView.

\begin{figure}[htbp]
  \centering
  \safeincludegraphics[width=1\linewidth]{psi_ref_ratio_diagnostic.png}
  \caption{Effect of the cortical--trabecular reference-stimulus ratio $r=\psi_{\mathrm{ref}}^{\mathrm{cort}}/\psi_{\mathrm{ref}}^{\mathrm{trab}}$ on similarity to the CT-derived population template (population reference).
  Left: nodal RMSE of the final density field (total and compartment-restricted, using CT-defined cortical/trabecular masks).
  Right: simulated cortical--trabecular contrast $C(r)$ compared to the CT contrast $C_{\mathrm{CT}}$.
  Generated by \texttt{analysis/psi\_ref\_ratio\_analysis.py}.}
  \label{fig:psi_ref_ratio_diagnostic}
\end{figure}

\begin{figure}[htbp]
  \centering
  \safeincludegraphics[width=0.7\linewidth]{psi_ref_ratio_histograms.png}
  \caption{Global density distributions at $T=500$ days for the CT template and selected ratios (default: $r=1$ and the best-RMSE run), illustrating how the ratio controls the overall density spectrum.
  Generated by \texttt{analysis/psi\_ref\_ratio\_analysis.py}.}
  \label{fig:psi_ref_ratio_histograms}
\end{figure}

\begin{figure}[htbp]
  \centering
  \safeincludegraphics[width=0.7\linewidth]{psi_ref_ratio_timecourse.png}
  \caption{Domain-mean density evolution $\bar\rho(t)$ for the ratio sweep, reconstructed from the conserved mass integral (\texttt{steps.csv}).
  The dashed line shows the CT template volume-mean density mapped to the FE mesh.
  Generated by \texttt{analysis/psi\_ref\_ratio\_analysis.py}.}
  \label{fig:psi_ref_ratio_timecourse}
\end{figure}
